\section{局限与未来工作}
\label{sec:limitations}

\paragraph{残差误差的性质.}
T5 的最大残差（近黑紫，$\Delta E_{00}\approx5.8$）与主套件的最大残差（大变明度差叠加色相对立的颜色对，10.7--23.8）同源：参考在四颜料基上表达其色域，红族与紫色颜料的高阶跨通道关联是该潜表示的性质，不是逐通道独立形式能够表示的。PBM 沿四个维度逼近潜空间行为——K--M 核心、幂律指数、谐波修正与 $\varphi$ 映射——但剩余误差仍主导最大值统计。

\paragraph{近中性深色盲区.}
近中性深色（$\mathrm{maxd}\in[0.02,0.06]$）落在门控之下，仅由核心混色。幂律阶梯是在有色度深色上拟合的，其通道结构与灰调深色所需的近似统一指数不同。残差在浅重侧测得最大值 9.5 $\Delta E_{00}$，T5E 全集中值 2.441。这位于设计目标之内，但在深色加白的设计目标套件（T5、T5E、tail）中是最大的最大值。$\varphi$ 映射减轻但未消除它。

\paragraph{与 Mixbox 的关系.}
Mixbox 采用潜空间线性混合加多项式解码，并在实测颜料光谱上训练。PBM 采用解析 K--M 结构与拟合自参考输出的解析修复层。由于参考的潜空间训练自实测颜料数据，PBM 传递性地继承参考的部分物理合理性；其相对真实颜料的误差被参考自身误差加上 PBM 相对参考的实测偏差所界定。该偏差是完全解析、无需训练表或许可限制的实现代价。

\paragraph{尚存的局限.}
\begin{enumerate}
\item 逐通道独立结构无法表示参考的高阶跨通道关联；这是最大误差的来源。
\item 常数拟合自 Mixbox~2.0。若参考改变，常数必须以 \cref{sec:estimation} 的流程重新估计；它们不能迁移到别的参考。
\item 线性方向响应无法用单一单变量形式同时协调品红与深红。T5 中品红成员的均值为 1.91--1.99、最大值为 2.19--2.86；近黑紫 $\#2A0A5E$ 均值 3.93、最大 5.82；深紫 $\#080248$ 为 5.14 / 7.33。
\item 谐波修正是手工特征模板上的线性模型，其残差在参考轨迹关于混合比非单调的颜色对上最大。非线性但仍解析的修正，以及模型侧新增潜通道，留待未来工作。
\item 合成光谱验证确认了模型在与独立两常数 K--M 世界中的平价性与定性族行为。实测量版本需要实测颜料光谱的授权发布。
\end{enumerate}

未来工作包括：扩展当前模板之外的特征结构、为近中性盲区重拟合阶梯、在核心中加入残差（非光谱）潜通道、把模型推广到分层合成，以及在可获得实测颜料数据时按实物重拟合。
